% Notes on endings with minor pieces. See README.md for the format.

% Diagram 78 of the 1500-1599 book: Mohamed - Alagouri
@ 8/8/8/4k1pN/8/1K2b3/8/8 b - - 1 81 | Kf5?
This is quite a difficult position. Unlike rook endings, there is no clear theory to rely on.
Both sides should know the typical tactics and use the \term{process-of-elimination} to avoid
mistakes.\index{Practical play!Process of elimination}

A few guidelines help. The defending king is far away, so the defender must not let his knight
lose its mobility.\index{Bishop against knight!Mobility of the knight} The defender can also offer to give up the knight
for the pawn: this forces the pawn forward, and then the knight can fork.\index{Bishop against knight!Knight fork}

81...\B f4! restricts the knight. 82.\K c3 \K f5 83.\N g7+ \K e4 84.\N h5 \B e5+, and the knight
cannot move at all. This is called \vocab[strategy]{domination}{domination}{A piece is attacked
and has no safe square to go to.}.\index{Bishop against knight!Domination} Black brings his king to the knight
and wins.

@ 8/8/8/6p1/5k2/1K2b1N1/8/8 w - - 4 83 | Nf1?
83.\N h5+! draws: 83...\K g4 84.\N g7 \K g3 85.\N h5+ \K h4 (85...\K g4 86.\N g7 (86.\N f6+?
\K f5 87.\N d5 \B d4 loses the knight) 86...\K f3 87.\K c4 g4 88.\N h5 \B f2 89.\K d5 \B h4)
86.\N f6 \K g3 (86...\B d4 87.\N e4 g4 88.\K c2 \B a7 89.\K d3).

@ 8/8/8/6p1/5k2/1K2b3/8/5N2 b - - 5 83 | Kf3?
83...\K e4! wins.

@ 8/8/8/6p1/8/1K2bk2/8/5N2 w - - 6 84 | Kc3?
84.\N h2+! draws: 84...\K g3 85.\N f1+ \K f2 86.\N h2 \K g2 87.\N g4 \K f3 88.\N h6 \B d4
89.\K c2 \B g7 90.\N f5 \B f6 91.\N h6 \B e7 92.\N f7 g4 93.\N e5+. The defender offers the
knight for the pawn, forces the pawn forward, and then forks.\index{Bishop against knight!Knight fork}

@ 8/8/8/6p1/8/2K1bk2/8/5N2 b - - 7 84 | Bf4!
\term{domination} again.\index{Bishop against knight!Domination}
